% ============================================================
% Section 183: 一维无限深方势阱的电偶极跃迁选择定则
% ============================================================
\section{量子力学：一维无限深方势阱的电偶极跃迁选择定则}
\label{sec:1d-well-dipole-selection}

\begin{modelbox}[题目：一维无限深方势阱的电偶极跃迁选择定则]
\begin{enumerate}
    \item 质量为 $m$ 的粒子束缚在宽为 $l$ 的一维无限深方势阱中，求基态和前两个激发态的能量，画出相应波函数；
    \item 计算从前两个激发态到基态间电偶极跃迁的矩阵元，解释定量差别；
    \item 给出系统任意两个态之间电偶极跃迁的一般选择定则。
\end{enumerate}
\end{modelbox}

\begin{tipbox}[考点快速分析]
\begin{itemize}
    \item \textbf{对称势阱}：$\psi_n^+(x)=\sqrt{2/l}\cos(n\pi x/l)$（$n$ 奇），$\psi_n^-(x)=\sqrt{2/l}\sin(n\pi x/l)$（$n$ 偶）
    \item \textbf{偶极算符}：$x$ 为奇宇称，$\langle f|x|i\rangle$ 非零要求初末态宇称相反
    \item \textbf{$2\to1$}：奇$\to$偶，允许；\textbf{$3\to1$}：偶$\to$偶，禁戒
    \item \textbf{一般定则}：$\Delta n$ 为奇数
\end{itemize}
\end{tipbox}

\begin{derivationbox}[标准解题过程]
\textbf{1. 能级与波函数}

对称势阱 $-l/2<x<l/2$：
\[
E_n=\frac{n^2\pi^2\hbar^2}{2ml^2},\quad n=1,2,3,\ldots
\]

\textit{基态} ($n=1$, 偶)：$\psi_1(x)=\sqrt{2/l}\cos(\pi x/l)$，无节点。

\textit{第一激发态} ($n=2$, 奇)：$\psi_2(x)=\sqrt{2/l}\sin(2\pi x/l)$，一个节点在 $x=0$。

\textit{第二激发态} ($n=3$, 偶)：$\psi_3(x)=\sqrt{2/l}\cos(3\pi x/l)$，两个节点。

\textbf{2. 电偶极跃迁矩阵元}

矩阵元 $\langle x\rangle_{fi}=\int_{-l/2}^{l/2}\psi_f^*(x)\,x\,\psi_i(x)\,dx$。

\textit{$2\to1$}：
\[
\langle x\rangle_{21}=\frac{2}{l}\int_{-l/2}^{l/2}x\cos\frac{\pi x}{l}\sin\frac{2\pi x}{l}\,dx.
\]
被积函数：$x$（奇）$\times\cos$（偶）$\times\sin$（奇）$=$ 偶。积分非零，跃迁允许。

\textit{$3\to1$}：
\[
\langle x\rangle_{31}=\frac{2}{l}\int_{-l/2}^{l/2}x\cos\frac{\pi x}{l}\cos\frac{3\pi x}{l}\,dx.
\]
被积函数：$x$（奇）$\times\cos$（偶）$\times\cos$（偶）$=$ 奇。对称区间积分为零，跃迁禁戒。

\textbf{3. 一般选择定则}

坐标算符 $x$ 为奇宇称。对称势阱中 $n$ 奇$\to$偶宇称，$n$ 偶$\to$奇宇称。因此 $\langle n'|x|n\rangle$ 非零要求
\begin{equation}
\boxed{\Delta n=n'-n\text{ 为奇数}.}
\end{equation}

即只有初末态量子数奇偶性不同时，电偶极跃迁才被允许。这是宇称守恒在跃迁过程中的直接体现。
\end{derivationbox}

\begin{formulabox}[答案汇总]
\begin{center}
\begin{tabular}{ll}
\toprule
\textbf{结论} & \textbf{结果} \\
\midrule
能级 & $E_n=n^2\pi^2\hbar^2/(2ml^2)$ \\
$2\to1$ 矩阵元 & 非零（宇称相反），跃迁允许 \\
$3\to1$ 矩阵元 & 零（宇称相同），跃迁禁戒 \\
一般选择定则 & $\Delta n$ 为奇数 \\
\bottomrule
\end{tabular}
\end{center}
\end{formulabox}

\begin{insightbox}[物理图像]
\begin{itemize}
    \item 一维势阱的电偶极选择定则 $\Delta n=$ 奇数是宇称守恒的直接结果。在更高维或更复杂的势中，选择定则由系统的对称群（如 $SO(3)$）决定。
    \item $3\to1$ 跃迁虽被电偶极禁戒，但可通过多极跃迁（如磁偶极、电四极）或碰撞诱导实现，只是概率远小于允许跃迁。
    \item 该模型是理解半导体量子阱光学吸收和发射光谱的起点：允许跃迁对应光吸收峰，禁戒跃迁在理想对称结构中不出现。
\end{itemize}
\end{insightbox}
